% =====================================================================
% EdgeFlow 用户手册 · 第 8 章 安全与认证
% 事实依据：docs/SECURITY.md、docs/API-SPEC.md、cloud/pkg/auth、cloud/pkg/audit、
%           pkg/certs、cmd/cloudcore/main.go、cmd/keadm
% =====================================================================
\chapter{安全与认证}
\label{ch:8}

EdgeFlow v0.1.0 的安全基线由三层构成：\textbf{云边通道 mTLS}（双向认证）、
\textbf{管理 API Token 认证}（默认关闭）与\textbf{审计台账}（默认开启）。
本章按层说明启用方式与运维要点。

\section{云边通道 mTLS}

云边 WebSocket 通道默认是明文 \texttt{ws://}（仅限可信网络/本机验证）。
生产环境必须启用 \textbf{mTLS}（双向 TLS）：云端验证边缘证书、边缘验证
云端证书与主机名，任何一侧校验失败连接在协议层之前被拒绝。

\subsection{启用开关}

\begin{table}[htbp]
\centering
\caption{mTLS 启用开关}
\label{tab:8-tls}
\begin{tabularx}{\textwidth}{l l X}
\toprule
\textbf{组件} & \textbf{开关（=on）} & \textbf{证书目录变量（默认）} \\
\midrule
云端 CloudHub & \texttt{EDGEFLOW\_CLOUDCORE\_TLS} & \texttt{EDGEFLOW\_CLOUDCORE\_CERT\_DIR}（\texttt{data/certs/}） \\
边缘 EdgeHub & \texttt{EDGEFLOW\_EDGECORE\_TLS} & \texttt{EDGEFLOW\_EDGECORE\_CERT\_DIR}（\texttt{data/certs/}） \\
\bottomrule
\end{tabularx}
\end{table}

\textbf{两侧必须同时开启}：云端只开 TLS 会拒绝明文连接；边缘只开 TLS
无法通过明文云端的握手。边缘开启后连接地址自动归一化为 \texttt{wss://}。

\begin{lstlisting}
# 云端
export EDGEFLOW_CLOUDCORE_TLS=on
export EDGEFLOW_CLOUDCORE_CERT_DIR=/data/certs
# 跨主机/集群部署必须注入 SAN（逗号分隔），非法条目启动即报错（fail-fast）
export EDGEFLOW_CLOUDCORE_TLS_SAN="IP:10.0.0.5,DNS:cloudcore.edgeflow.svc"
bin/cloudcore

# 边缘
export EDGEFLOW_EDGECORE_TLS=on
export EDGEFLOW_EDGECORE_CERT_DIR=/data/certs
export EDGEFLOW_EDGECORE_CLOUD_ADDR=wss://10.0.0.5:10000/v1/edge
bin/edgecore
\end{lstlisting}

\textbf{关键约束}：边缘连接的地址（\texttt{wss://<host>:<port>}）必须落在
服务端证书 SAN 内，否则握手失败。未设置 SAN 时证书仅覆盖本机回环
（\texttt{127.0.0.1/localhost/cloudcore}），mTLS 默认只在本机可用。

\subsection{使用 keadm 生成 TLS 产物}

\begin{lstlisting}
# 云端产物：注入 TLS 开关 + 证书目录 + SAN
keadm init --tls --tls-san=IP:192.168.1.10 --output-dir=./keadm-out

# 边缘产物：注入 EDGEFLOW_EDGECORE_TLS=on + CERT_DIR，地址用 wss://
keadm join --cloudcore-ip=192.168.1.10 --cloudcore-port=31000 \
  --token=<token> --node-id=edge-worker-01 --tls --output-dir=./keadm-out
\end{lstlisting}

\subsection{设备接入令牌认证}

边缘节点注册接入的令牌校验（设备认证）
\textbf{已实现}，链路如下：

\begin{itemize}
  \item \texttt{keadm join} 生成 \texttt{edgecore.env} 时写入
        \texttt{EDGEFLOW\_EDGECORE\_TOKEN}（值为 \texttt{--token} 参数）；
  \item edgecore 启动后随 \texttt{Register} 消息\textbf{携带}该令牌上报云端；
  \item 云端设置 \texttt{EDGEFLOW\_CLOUDCORE\_NODE\_TOKEN} 后启用校验：
        按\textbf{常数时间比较}，不匹配 → \textbf{拒绝注册}；
  \item 云端未设置该变量 → 不校验，保持向后兼容（默认空 = 不校验）。
\end{itemize}

\begin{lstlisting}
# 云端：设置后开启设备接入校验（建议生产启用）
export EDGEFLOW_CLOUDCORE_NODE_TOKEN=<secret>
bin/cloudcore
\end{lstlisting}

\section{证书管理}

\subsection{证书布局}

证书目录（云端与边缘各自 \texttt{data/certs/}）包含：

\begin{table}[htbp]
\centering
\caption{证书文件布局}
\label{tab:8-certs}
\begin{tabularx}{\textwidth}{l l l X}
\toprule
\textbf{文件} & \textbf{用途} & \textbf{有效期} & \textbf{关键属性} \\
\midrule
\texttt{ca.crt} / \texttt{ca.key} & 自签 CA & 10 年 & \texttt{CN=edgeflow-ca}，CertSign \\
\texttt{cloudcore.crt} / \texttt{cloudcore.key} & 云服务端证书 & 1 年 & SAN：回环 + 注入值，EKU serverAuth \\
\texttt{edgecore.crt} / \texttt{edgecore.key} & 边客户端证书 & 1 年 & \texttt{CN=edgeflow-<nodeID>}，EKU clientAuth \\
\texttt{crl.pem} & 吊销列表（X.509 CRL 签名产物） & 7 天（有新吊销时重签） & mTLS 握手校验依据 \\
\texttt{crl.json} & 吊销序列号来源记录 & — & 与 \texttt{crl.pem} 对账自愈，唯一事实源 \\
\texttt{crl.lock} & 吊销写路径文件锁 & — & 写路径互斥（flock） \\
\bottomrule
\end{tabularx}
\end{table}

私钥权限 \texttt{0600}，证书文件 \texttt{0644}。证书在组件首次以 TLS
启动时\textbf{自动生成}（幂等：已存在则加载校验，绝不重新生成）；也可用
\texttt{hack/gen-certs.sh} 预置。证书损坏/不完整 → \textbf{启动失败}
（fail-fast，不降级）。

\subsection{证书轮换与吊销}

\textbf{人工编排路径}：

\begin{enumerate}
  \item 签发新证书（先备份并删除旧文件，或使用新证书目录）；
  \item 滚动重启组件，\textbf{先云后边}：云端加载新服务端证书，边缘
        重连时自动携带新客户端证书；
  \item 验证新连接建立后，清除旧证书文件。
\end{enumerate}

轮换 CA 需\textbf{全量重签}所有叶子证书（旧叶子证书随 CA 更换立即失效）。
CA 私钥仅靠文件权限保护，生产建议离线签发后移出节点或接入 KMS。

\textbf{自动轮换（\texttt{keadm cert rotate}）}：
\texttt{keadm cert rotate --node=<CN> --cert-dir=<目录>} 按 CN 定位证书并
\textbf{事务化重签}（SAN 继承），旧证书备份到 \texttt{backups/<时间戳>/}
（含 \texttt{manifest.json}），输出新证书路径；随后需将新证书
\textbf{重新分发}到目标节点并重启生效。

\textbf{吊销（\texttt{keadm cert revoke}）}：
\texttt{keadm cert revoke --node=<CN>|--serial=<hex> --cert-dir=<目录>}

\begin{itemize}
  \item \texttt{--node} 与 \texttt{--serial} \textbf{二选一互斥}：
        \texttt{--node} 按 CN 定位取序列号；\texttt{--serial} 接受十六进制
        （可带 \texttt{0x} 前缀，大小写均可，自动归一化），\textbf{不依赖
        证书文件}——可直接吊销已轮换的历史序列号（防泄露）；
  \item 序列号写入 \texttt{crl.json} 并重签 \texttt{crl.pem}；重复吊销
        \textbf{幂等}（无副作用）；
  \item 吊销\textbf{不删除证书文件}：已分发的证书由对端在 mTLS 握手时
        按 CRL \textbf{拒绝}（cloudcore/edgecore 双侧消费端均接线，真实
        握手测试覆盖）；
  \item 吊销后需 \texttt{keadm cert rotate} 重签新证书并\textbf{重新分发}。
\end{itemize}

\textbf{吊销语义要点}：

\begin{itemize}
  \item CRL 缺失（无 \texttt{crl.pem}）视为\textbf{无吊销}（向后兼容放行）；
  \item CRL 仅在发生新吊销时重生成（\texttt{NextUpdate} = 生成时 +7 天）；
        过期后\textbf{已吊销序列号仍被拒绝}（吊销检查先于过期检查），未吊销
        证书默认\textbf{放行}；\texttt{FailOnExpired} 为库级选项，当前
        \textbf{未接线}到 mTLS 握手；
  \item \texttt{crl.json}/\texttt{crl.pem} 损坏或签名无效 → \textbf{拒绝}
        （fail-closed，不静默放行）；
  \item 吊销写路径文件锁 \texttt{crl.lock} 获取失败时自动\textbf{降级为无锁
        校验}（功能语义不变，仅损失 \texttt{crl.json} 领先时的即时重生成
        自愈），并输出 \textbf{5 分钟限频}的 Warn 告警日志，便于发现证书
        目录权限异常；
  \item 吊销\textbf{只增不减}：误吊销后无撤销命令，只能重新签发新证书
        （新序列号）；\textbf{请勿手工编辑 crl.json}（下次吊销时会按记录
        精确重生成 CRL，手工删除的条目将材料化生效）。
\end{itemize}

\subsection{在线吊销查询（OCSP）}

云端提供标准 OCSP responder（RFC 6960）：\texttt{POST /ocsp}。

\begin{itemize}
  \item 请求体：DER 编码 \texttt{OCSPRequest}（\texttt{Content-Type:
        application/ocsp-request}）；响应体：DER 编码 \texttt{OCSPResponse}
        （\texttt{application/ocsp-response}），由\textbf{CA 私钥签名}；
  \item 吊销状态直接读取 \texttt{crl.json}，与 CRL \textbf{同源}——吊销后
        双通道（CRL 握手校验 + OCSP 在线查询）同时生效；
  \item 免认证端点，以 per-IP 令牌桶限流防滥用：默认 10 req/s、
        burst 20，超限返回 \texttt{429}；速率可经环境变量
        \texttt{EDGEFLOW\_CLOUDCORE\_OCSP\_RATE\_LIMIT} 调整；成功响应带
        \texttt{Cache-Control: max-age=3600}（\texttt{nextUpdate} ≈ 7 天）；
  \item 状态码：\texttt{200} = 签名响应；\texttt{400} = 请求非法/超限
        （请求体上限 \textbf{16\,KiB}，空请求/DER 损坏/超限）；\texttt{429} =
        触发 per-IP 限流；\texttt{500} =
        CA 不可用/吊销记录损坏/响应构建失败（fail-closed，不静默放行）；
  \item 客户端查询库函数 \texttt{OCSPStatus/OCSPStatusAt} 可用（验签 +
        responderID + CertID 匹配），新增新鲜度校验入口
        \texttt{OCSPStatusAtWithPolicy}/\texttt{ParseOCSPResponseWithFreshness}
        （fail-closed：过期/未来时间拒绝，默认 5 分钟 skew）；当前均
        \textbf{未接入}任何命令或握手路径——mTLS 握手消费的是 CRL 而非
        OCSP（生产路径接入时须用 WithPolicy 入口）。
\end{itemize}

\textbf{安全边界}：\texttt{/ocsp} 是唯一\textbf{免 Token 认证}的协议端点
（响应由 CA 私钥签名，客户端必须验签防伪造）。该端点已内置 per-IP 限流
（默认 10 req/s、burst 20，超限 429），但仍监听全部接口，建议置于受信
内网/防火墙之后（per-IP 粒度对分布式多 IP 放大不设防，生产可叠加 LB 层
全局限流）；吊销状态属公开信息（协议设计如此），不涉及凭据泄露。

\section{管理 API Token 认证}

管理 API（REST，8080 端口）默认\textbf{不认证}（向后兼容）。启用方式：

\begin{lstlisting}
export EDGEFLOW_CLOUDCORE_AUTH=on
export EDGEFLOW_CLOUDCORE_API_TOKEN=<你的令牌>   # 必填，未设置则启动失败
bin/cloudcore
\end{lstlisting}

\begin{itemize}
  \item \texttt{EDGEFLOW\_CLOUDCORE\_AUTH=on} 且 \texttt{API\_TOKEN} 未设置
        → \textbf{启动失败}（fail-fast，防止"开了认证却没令牌"的空转）；
  \item 请求必须携带 \texttt{Authorization: Bearer <token>}；
  \item 缺失/非 Bearer 方案/令牌不匹配 → \texttt{401} +
        \texttt{WWW-Authenticate: Bearer} 响应头；
  \item 常量时间比较防时序侧信道；认证通过者记为身份 \texttt{token}
        （单令牌模型，持有令牌即管理员；完整 RBAC 角色模型为后续版本）。
\end{itemize}

\section{审计台账}

管理 API 的所有操作默认写入审计台账：

\begin{itemize}
  \item 路径：\texttt{EDGEFLOW\_CLOUDCORE\_AUDIT\_PATH}，默认
        \texttt{data/audit-ledger.jsonl}；
  \item 格式：\textbf{JSONL}（逐行 JSON，追加写，不覆盖历史）；
  \item 记录内容：请求路径、方法、结果、身份（Token 认证下含
        operator；认证失败记录 \texttt{anonymous}）；
  \item \textbf{审计失败不阻断 API}：写盘失败只记错误日志，业务请求照常
        处理（可用性优先）。
\end{itemize}

\begin{lstlisting}
# 查看审计台账（每行一条记录）
tail -f data/audit-ledger.jsonl
\end{lstlisting}

\section{安全基线小结}

\begin{itemize}
  \item 已实现：双向认证（云验边、边验云+主机名）、设备接入令牌认证
        （云端设 \texttt{EDGEFLOW\_CLOUDCORE\_NODE\_TOKEN} 后启用，见
        「设备接入令牌认证」小节）、私钥 0600、强制 TLS 1.2+、证书
        损坏 fail-fast、Token 认证（默认关）、审计台账（默认开）、
        证书吊销（CRL 离线 + OCSP 在线，见「证书轮换与吊销」与
        「在线吊销查询（OCSP）」小节）；
  \item 已知限制：CSR 审批流未实现；证书 CN 与 nodeID 未绑定校验；
        单令牌模型（无多角色授权）。
\end{itemize}
